\section{CORRENTI STAZIONARIE}

\subsection{Densit\`a di corrente e conservazione della carica}

Quando fra i capi di un filo conduttore viene mantenuta, per esempio tramite una
pila, una differenza di potenziale, le cariche libere all'interno del conduttore
non possono essere in equilibrio (il conduttore non \`e equipotenziale), e si ha
una \textbf{corrente elettrica} nel filo conduttore, cio\`e un moto di cariche.
Fissato un orientamento sul filo, si definisce \textbf{intensit\`a di corrente}
$I$ la quantit\`a di carica che attraversa una sezione $S$ del filo nell'unit\`a
di tempo nella direzione presa come verso positivo:
\begin{equation}
  I(t)\defeq\frac{dQ}{dt}
  \tag{3.1}
\end{equation}
dove $Q(t)$ \`e la quantit\`a di carica che attraversa $S$ nell'intervallo di
tempo $(t_0,\ t)$ ($t_0$ \`e arbitrario). In generale $I$ dipende sia dal tempo
(se la differenza di potenziale non \`e costante), che dalla posizione della
sezione $S$. L'intensit\`a di corrente pu\`o essere positiva o negativa: per
esempio pu\`o essere negativa o perch\'e $S$ \`e attraversata da cariche
negative nel verso assunto come positivo, oppure perch\'e cariche positive
attraversano $S$ nel verso negativo.\\
Nel SI l'unit\`a di misura di corrente \`e l'amp\`ere (\texttt{A}): si ha una
corrente di 1\,amp\`ere attraverso una superficie $S$ quando la carica di
1\,coulomb attraversa $S$ in 1\,s. Dalla (1.3) si ha
\begin{equation}
  1\,\mathrm A=3\times10^9\,\mathrm{ues}.
  \tag{3.2}
\end{equation}
\`E ovvio che si pu\`o avere un moto di cariche, e quindi una corrente, non solo
in un filo conduttore, ma pi\`u in generale in un conduttore di forma qualsiasi
e anche nel vuoto, come per esempio nel tubo di un televisore dove gli elettroni
estratti per effetto termoionico da un filamento caldo vengono accelerati da
un'opportuna differenza di potenziale.\\
Possiamo pensare alle cariche in moto come ad un fluido, e definire la
\textbf{densit\`a di corrente} $\vec J(x,y,z,t)$: supponiamo di avere una
densit\`a di carica $\rho(x,y,z,t)$ e sia $\vec u(x,y,z,t)$ il campo di
velocit\`a delle cariche: $\vec u(x,y,z,t)$ \`e la velocit\`a (media) delle
cariche che all'istante $t$ si trovano in un intorno del punto $(x,y,z)$. Allora
\begin{equation}
  \vec J(x,y,z,t)\defeq\rho(x,y,z,t)\vec u(x,y,z,t).
  \tag{3.3}
\end{equation}
Pi\`u in generale, la corrente potrebbe essere dovuta a pi\`u distribuzioni,
ciascuna con il proprio campo di velocit\`a, come per esempio in un gas o in un
liquido dove si hanno portatori di carica (ioni) sia positivi che negativi che
si muovono con velocit\`a diverse, o anche in certi acceleratori di particelle
dove vengono fatti circolare in senso inverso sia elettroni che positroni
(elettroni positivi). In questi casi la (3.3) si generalizza nel modo seguente:
\begin{equation}
  \vec J(x,y,z,t)=\textstyle\sum_i\rho_i(x,y,z,t)\vec u_i(x,y,z,t)
  \tag{3.3'}
\end{equation}
cio\`e la densit\`a di corrente \`e la somma delle densit\`a di corrente
associate ai diversi portatori di carica in moto con velocit\`a diverse: in un
conduttore in quiete abbiamo
gli elettroni di conduzione con un certo campo di velocit\`a, e gli ioni del
reticolo con densit\`a opposta ma campo di velocit\`a nullo.\\
Grazie al significato di flusso (v.\ capitolo 1.4), la quantit\`a di carica che
attraversa nell'unit\`a di tempo una superficie orientata $S$, cio\`e la
corrente $I(S)$, \`e il flusso di $\vec J$ attraverso $S$:
\begin{equation}
  I(S)=\int_S\vec J\cdot d\vec S.
  \tag{3.4}
\end{equation}

\paragraph{Esempio}
\begin{quote}
  In un filo conduttore le cariche in moto sono solo gli elettroni, per cui la
  densit\`a di corrente \`e data dalla (3.3). Vogliamo calcolare la velocit\`a
  che compare nella (3.3). Supponiamo che il filo abbia una sezione $S$ di
  1\,mm$^2$ e sia percorso da una corrente di 1\,A. Siccome in un conduttore ci
  sono circa $10^{23}$\,elettroni di conduzione/cm$^3$, si ha
  \begin{equation}
    \begin{aligned}
      &\rho=1.6\times10^{-19}\cdot10^{23}=1.6\times10^4\,\mathrm{C/cm^3}
      \quad\Rightarrow\\
      &u=\frac{I}{\rho S}=\frac{1}{10^{-2}\cdot1.6\times10^4}\,\mathrm{cm/s}
      \approx10^{-2}\,\mathrm{cm/s}.
    \end{aligned}
    \tag{3.5}
  \end{equation}
\end{quote}
Un valore cos\`i piccolo per la velocit\`a degli elettroni ci fa capire che si
tratta di una velocit\`a media: gli elettroni di conduzione costituiscono un gas
e, come le particelle di ogni gas, sono animati da un moto disordinato con
energia cinetica media dell'ordine di $kT$ e quindi, a temperatura ambiente,
velocit\`a termica dell'ordine di $\sqrt{kT/m}\simeq10^7$\,cm/s. Questa che
abbiamo appena calcolato \`e la ``velocit\`a quadratica media'', cio\`e
$(\overline{v^2}\,)^{1/2}$: in condizioni normali qualsiasi componente della
velocit\`a ha media nulla, in quanto non c'\`e nessuna direzione privilegiata
(distribuzione isotropa delle velocit\`a); in presenza di una differenza di
potenziale ai capi del filo, e quindi di un campo elettrico all'interno del
filo, a questa distribuzione isotropa delle velocit\`a si sovrappone una piccola
velocit\`a d'insieme nella direzione del filo, il cui valore medio \`e dato
dalla (3.5).

Una delle leggi pi\`u importanti e meglio verificata della fisica \`e la legge
di conservazione della carica elettrica: la carica totale di un sistema isolato
resta costante nel tempo. Quando negli esperimenti di fisica delle particelle
elementari vengono prodotte particelle cariche, queste vengono sempre prodotte
insieme ad altre particelle in modo tale che la carica totale delle particelle
prodotte \`e nulla: per esempio nell'urto di raggi $\gamma$ (fotoni di altissima
energia) con nuclei pesanti vengono prodotte coppie di elettroni e positroni:
$\gamma\to e^++e^-$, ed \`e sperimentalmente ben verificato che la carica del
positrone \`e opposta a quella dell'elettrone. Un altro esempio \`e costituito
dal decadimento del neutrone \ $n\to p+e^-+\overline\nu$: \ insieme
all'elettrone viene prodotto un protone ed un antineutrino; sperimentalmente \`e
stato stabilito che la carica del neutrone e quella del neutrino sono minori di
$10^{-21}e$, e con la stessa precisione che \`e nulla la carica del protone
pi\`u quella dell'elettrone: quindi la legge di conservazione della carica \`e
stabilita con una precisione dell'ordine di una parte su $10^{21}$.\\
La legge di conservazione della carica si pu\`o esprimere in questo modo: se $S$
\`e la superficie che racchiude un volume $V$, e $Q_V(t)$ \`e la carica
contenuta in esso all'istante $t$, la corrente uscente da $S$ \`e uguale alla
diminuzione di $Q_V(t)$ nell'unit\`a di tempo: se
attraverso $S$ esce della carica, si ha una corrispondente diminuzione della
carica contenuta in $V$, e viceversa (ci\`o non sarebbe vero se dentro $V$ le
cariche nascessero dal nulla o scomparissero nel nulla). Formalmente:
\begin{equation}
  \Phi_S(\vec J)=-\frac{dQ_V(t)}{dt}\,\cdot
  \tag{3.6}
\end{equation}
La (3.6) pu\`o essere riscritta in forma \underline{locale}, nello stesso senso
in cui il teorema di Gauss \`e tradotto in forma locale dall'equazione
$\dive\vec E=4\pi\rho$: dal teorema della divergenza
\begin{equation}
  \Phi_S(\vec J)\equiv\int_S\vec J\cdot d\vec S=\int_V\dive\vec J(x,y,z,t)\,dV
  \tag{3.7}
\end{equation}
e poich\'e
\begin{equation}
  \frac{dQ_V(t)}{dt}=\frac{d}{dt}\int_V\rho(x,y,z,t)\,dV
  =\int_V\frac{\partial\rho(x,y,z,t)}{\partial t}\,dV,
  \tag{3.8}
\end{equation}
(\`e lecito portare la derivata rispetto a $t$ dentro l'integrale), dalla (3.6)
si ha
\begin{equation}
  \int_V\dive\vec J(x,y,z,t)\,dV
  =-\int_V\frac{\partial\rho(x,y,z,t)}{\partial t}\,dV
  \tag{3.9}
\end{equation}
e per l'arbitrariet\`a del volume $V$
\begin{equation}
  \dive\vec J+\frac{\partial\rho}{\partial t}=0
  \tag{3.10}
\end{equation}
nota come \textbf{equazione di conservazione della carica}: le sorgenti di
$\vec J$ sono i punti in cui si ha una diminuzione di carica
($\partial\rho/\partial t<0$), i pozzi quelli in cui si ha un accumulo.

In condizioni stazionarie (per definizione) $\vec J$ e $\rho$ sono indipendenti
dal tempo, quindi dalla (3.10) per una \textbf{corrente stazionaria}
\begin{equation}
  \dive\vec J(x,y,z)=0
  \tag{3.11}
\end{equation}
da cui segue che la corrente uscente da una qualsiasi superficie chiusa \`e
zero. In particolare, se la corrente \`e confinata dentro un filo conduttore, il
filo stesso \`e un tubo di flusso, e quindi la corrente che attraversa una
sezione del filo \`e indipendente dalla posizione della sezione, dal suo
orientamento e dal tempo: si dice che il filo (o il circuito di cui il filo fa
parte) \`e percorso da una \textbf{corrente continua} $I$; per questa ragione la
(3.11) \`e detta \textbf{equazione di continuit\`a} \big(talvolta useremo questo
termine anche per l'equazione (3.10)\big).\\
Nel seguito avremo a che fare con le cosiddette ``correnti lentamente
variabili'': per esempio sono correnti lentamente variabili quelle che scorrono
nei circuiti alimentati dalla forza elettromotrice a 50\,hertz di cui disponiamo
nelle nostre case. La densit\`a di una corrente lentamente variabile \`e della
forma
\begin{equation}
  \vec J(x,y,z,t)=\vec J_0(x,y,z)f(t)
  \tag{3.12}
\end{equation}
con $f(t)$ una funzione lentamente variabile, un concetto che chiariremo nel
seguito; \underline{se}, come spesso accade (circuiti senza condensatori),
\ $\dive\vec J_0(x,y,z)=0$ \ allora anche le correnti lentamente variabili sono
a divergenza nulla per cui la corrente $I(t)$ in un circuito, pur dipendendo dal
tempo, \`e la stessa attraverso qualsiasi sezione del circuito; quella a
50\,hertz delle nostre case \ $I(t)=I_0\cos\omega t$ \ si chiama ``corrente
alternata''.

\subsection{La legge di Ohm}

Se fra due superfici di un conduttore (p.es.\ fra i capi di un filo) si ha una
differenza di potenziale $V$, la corrente $I$ nel conduttore \`e proporzionale a
$V$:
\begin{equation}
  V=RI
  \tag{3.13}
\end{equation}
dove $R$, il coefficiente di proporzionalit\`a, \`e detta la \textbf{resistenza}
del conduttore. La (3.13), nota come la \textbf{legge di Ohm}, \`e stata
stabilita da G.S. Ohm nel 1827: non \`e una legge fondamentale
dell'elettromagnetismo, ma piuttosto segue da specifiche propriet\`a della
materia in cui si muovono i portatori di carica. Quindi, come discuteremo, anche
la sua validit\`a non ha carattere generale, ma costituisce un'ottima
approssimazione in molte situazioni sperimentali.\\
Nel SI l'unit\`a di misura di resistenza \`e l'ohm \ $(\Omega)$: un conduttore
ha la resistenza di 1\,ohm quando una differenza di potenziale di 1\,volt
produce una corrente di 1\,amp\`ere. Nel sistema di Gauss la resistenza ha le
dimensioni di un tempo diviso una lunghezza, cio\`e l'inverso di una velocit\`a,
e
\begin{equation}
  1\,\Omega=\frac{1\,\mathrm V}{1\,\mathrm A}
  =\frac{1}{300\cdot3\times10^9}\,\mathrm{ues}
  =1.1\times10^{-12}\,\mathrm{ues}\,;
  \quad\ 1\,\mathrm{ues}=9\times10^{11}\,\Omega.
  \tag{3.14}
\end{equation}
Nella figura 3.1 \`e riportato il collegamento in serie ed in parallelo delle
resistenze con il simbolo per le resistenze usato nei circuiti elettrici. Le
resistenze in serie si sommano: per la (3.11) la corrente $I$ nelle due
resistenze \`e la stessa, quindi $V_1=R_1I$, \ $V_2=R_2I$, \ da cui
$V=V_1+V_2=(R_1+R_2)I$, quindi $R=R_1+R_2$. Se le resistenze sono collegate in
parallelo si sommano gli inversi delle resistenze: poich\'e $V=R_1I_1$ e
$V=R_2I_2$ e, sempre per la (3.11), la corrente totale \`e $I=I_1+I_2$, \ si ha
\ $1/R\equiv I/V=1/R_1+1/R_2$.
\fig{fig3_1}{3.1}
Anche in questo caso conviene esprimere la legge di Ohm in forma locale. Innanzi
tutto osserviamo che la resistenza \`e definita per una porzione di conduttore
che sia un tubo di flusso di $\vec J$ (perch\'e in un tubo di flusso \`e ben
definita la corrente che vi scorre), chiuso da due superfici equipotenziali.
Consideriamo all'interno di un conduttore un tubo di flusso infinitesimo di
$\vec J$: sia $\Delta S$ la sua sezione, $\Delta L$ la lunghezza ed $R$ la sua
resistenza. Se il conduttore \`e costituito da un materiale isotropo
(propriet\`a fisiche indipendenti dalla direzione), $\vec J$ ed $\vec E$ sono
paralleli e dalla (3.13) si ha
\begin{equation}
  V=E\,\Delta L=RI=R\,J\Delta S
  \quad\ \Rightarrow\quad\ \vec E=\frac{R\,\Delta S}{\Delta L}\vec J
  \tag{3.15}
\end{equation}
cio\`e definita la \textbf{resistivit\`a} del mezzo
\begin{equation}
  \varrho_c(x,y,z)\defeq\lim\frac{R\,\Delta S}{\Delta L}
  \tag{3.16}
\end{equation}
(il limite \`e per $\Delta L,\Delta S\to0$), si ha
\begin{equation}
  \begin{aligned}
    &\vec E(x,y,z)=\varrho_c(x,y,z)\vec J(x,y,z)\,;\\
    &\vec J(x,y,z)=\sigma_c(x,y,z)\vec E(x,y,z)
    \qquad\quad(\sigma_c\equiv\varrho_c^{-1}).
  \end{aligned}
  \tag{3.17}
\end{equation}
La resistivit\`a $\varrho_c$ \`e ben definita dalla (3.16) (cio\`e il limite di
$R\Delta S/\Delta L$ esiste) come conseguenza del fatto, contenuto nella
derivazione della (3.15), che per la legge delle resistenze in serie, se il
mezzo \`e omogeneo (e in un intorno infinitesimo di un punto lo \`e) $R$ \`e
proporzionale a $\Delta L$, mentre per quella delle resistenze in parallelo $R$
\`e inversamente proporzionale a $\Delta S$. La resistivit\`a \`e quindi una
propriet\`a (locale) del conduttore: in altri termini, mentre la resistenza di
un conduttore dipende non soltanto dal materiale (rame, argento \dots) di cui
\`e costituito, ma anche dalle sue dimensioni, dalla sua forma e da come viene
collegato al resto del circuito, $\varrho_c$ dipende solo dal tipo di materiale;
$\sigma_c$, che \`e semplicemente l'inverso di $\varrho_c$, \`e la
\textbf{conducibilit\`a} del mezzo.

La (3.17) ha gli stessi limiti di validit\`a della legge di Ohm, di cui \`e
l'espressione locale; in particolare il campo elettrico nel conduttore non deve
essere troppo intenso, altrimenti altera il meccanismo della conduzione, e tanto
$\varrho_c$ quanto $\sigma_c$ dipenderebbero dal campo elettrico: la (3.17)
costituisce l'\underline{approssimazione lineare} della (in generale complicata)
relazione esistente fra campo elettrico e densit\`a di corrente.\\
Se il mezzo che stiamo considerando \`e omogeneo, $\varrho_c$ \`e costante; in
questo caso, per un conduttore di sezione costante $S$ e lunghezza $L$ si ha
\begin{equation}
  R=\varrho_c\,\frac LS\,\cdot
  \tag{3.18}
\end{equation}
La (3.18) vale per correnti stazionarie o lentamente variabili: in questo caso
tutto il conduttore \`e un tubo di flusso e la densit\`a di corrente \`e
uniforme dentro il conduttore, mentre per correnti rapidamente variabili la
corrente scorre prevalentemente sulla superficie (``effetto pelle'') e quindi la
resistenza risulta maggiore di quella data dalla (3.18).\\
Dalla (3.16) si vede che le dimensioni di $\varrho_c$ sono quelle di una
resistenza per una lunghezza, e normalmente il suo valore viene espresso in
$\Omega\cdot$cm: per esempio, la resistivit\`a del rame \`e
$1.7\times10^{-6}\,\Omega\,$cm, e la resistenza di un filo di rame lungo 10\,m e
di sezione 1\,mm$^2$ per la (3.16) \`e
\ $R=1.7\times10^{-6}\cdot10^3/10^{-2}=0.17\,\Omega$. \ Nel sistema di Gauss la
resistivit\`a si misura in secondi; per la (3.14) la resistivit\`a del rame \`e
\ $1.9\times10^{-18}$\,ues.

\paragraph{Esempi}
\begin{enumerate}
  \item Vogliamo calcolare la resistenza di una corona sferica conduttrice di
    raggi $r_1$ e $r_2$, fra le sue superfici interna ed esterna.\\
    Consideriamo una corona sferica di raggio $r$ e spessore infinitesimo
    $\Delta r$; per la (3.18) la sua resistenza \`e
    \ $\varrho_c\Delta r/4\pi r^2$, quindi per la legge delle resistenze in
    serie
    \begin{equation}
      R=\varrho_c\int_{r_1}^{r_2}\frac{dr}{4\pi r^2}
      =\frac{\varrho_c}{4\pi}\Big(\frac{1}{r_1}-\frac{1}{r_2}\Big)\,\cdot
      \tag{3.19}
    \end{equation}
  \item La corrente circola (in senso rotatorio attorno all'asse) in un cilindro
    conduttore cavo di raggio interno $a$, spessore $s$ e altezza $h$. Vogliamo
    calcolare la resistenza del circuito.\\
    Questa volta il circuito \`e il \underline{parallelo} delle corone
    cilindriche di raggio $r$ \ $(a\le r\le a+s)$, spessore $\Delta r$ e altezza
    $h$; la lunghezza di ciascuno di questi circuiti elementari \`e \ $L=2\pi r$
    \ e la sezione \ $\Delta S=h\Delta r$, e per la (3.18) la resistenza \`e
    \ $2\pi\varrho_cr/h\Delta r$. Per la legge delle resistenze in parallelo,
    conviene calcolare l'inverso della resistenza:
    \begin{equation}
      \begin{aligned}
        &R^{-1}=(2\pi\varrho_c)^{-1}h\int_a^{a+s}\frac{dr}{r}
        =(2\pi\varrho_c)^{-1}h\log\frac{a+s}{a}\\
        &\qquad\Rightarrow R=\frac{2\pi\varrho_c}{h\log\big((a+s)/a\big)}\,\cdot
      \end{aligned}
      \tag{3.20}
    \end{equation}
    Se invece la corrente scorre fra la superficie interna e quella esterna del
    cilindro, analogamente alla (3.19) si ottiene
    \begin{equation}
      R=\varrho_c\int_a^{a+s}\frac{dr}{2\pi hr}
      =\frac{\varrho_c}{2\pi h}\log\frac{a+s}{a}
      \tag{3.21}
    \end{equation}
    diversa dalla (3.20), a conferma del fatto che la resistenza di un
    conduttore dipende da come viene collegato al resto del circuito.
\end{enumerate}

\subsection{Il lavoro del campo elettrico}

Un conduttore fra i cui estremi c'\`e una differenza di potenziale $V$, \`e
percorso da una corrente stazionaria $I$. Ci\`o significa che nell'intervallo di
tempo $\Delta t$ una carica pari a $I\Delta t$ entra nel conduttore dall'estremo
a potenziale maggiore e una identica quantit\`a di carica esce dall'altro
estremo (figura 3.2): ci\`o \`e equivalente al passaggio della carica
$I\Delta t$ dal potenziale $V$ al potenziale 0, quindi il campo elettrico ha
fatto sulle cariche del conduttore il lavoro $\mathcal L=VI\,\Delta t$, e quindi
una potenza (lavoro per unit\`a di tempo)
\fig{fig3_2}{3.2}
\begin{equation}
  W_\mathrm J=I\,V.
  \tag{3.22}
\end{equation}
Se sugli elettroni di conduzione non agissero altre forze, il lavoro fatto dal
campo elettrico produrrebbe un aumento della loro energia cinetica, cio\`e
un'accelerazione nella direzione del campo e quindi un aumento della corrente
(una differenza di potenziale di un volt accelererebbe un elettrone
\underline{libero} ad una velocit\`a di circa $10^5$\,m/s). Ma, a causa della
presenza del reticolo entro cui si muovono, e con il quale gli elettroni
interagiscono, su una scala di tempi macroscopica la loro energia cinetica non
cambia e continuano a muoversi nella direzione del campo elettrico con la loro
irrisoria velocit\`a media poich\'e ad ogni urto con gli ioni del reticolo
perdono l'energia cinetica acquisita tra un urto e l'altro: si tratta del
contatto termico fra due gas, il gas di elettroni ed il reticolo; il campo
elettrico fa lavoro sul gas di elettroni aumentandone l'energia, ma questa
energia rapidamente si equipartisce fra i due gas. Quindi aumenta l'energia
interna del conduttore, che si scalda. Questo effetto \`e noto come
\textbf{effetto Joule}, e la potenza (3.22) eguaglia la potenza dissipata dagli
elettroni ed \`e detta \textbf{potenza Joule}. Grazie alla (3.13) la (3.22)
pu\`o essere anche scritta come
\begin{equation}
  W_\mathrm J=I^2R=\frac{V^2}{R}\,\cdot
  \tag{3.22'}
\end{equation}
In generale, sia data una densit\`a di corrente $\vec J$, non necessariamente in
un conduttore, ed un campo elettrico $\vec E$. Sia $\Delta V$ un volumetto
infinitesimo, $\rho(x,y,z)\Delta V$ la carica in esso contenuta (ad un certo
istante), e $\vec u(x,y,z)$ la sua velocit\`a. Il lavoro che il campo elettrico
fa nell'unit\`a di tempo (la potenza) sulla carica in $\Delta V$ \`e
\begin{equation}
  W=\rho(x,y,z)\Delta V\,\vec E\cdot\vec u(x,y,z)=\vec E\cdot\vec J\,\Delta V
  \tag{3.23}
\end{equation}
e quindi la \textbf{potenza} $w$ \textbf{per unit\`a di volume} \`e
\begin{equation}
  w(x,y,z)=\vec E(x,y,z)\cdot\vec J(x,y,z)\,.
  \tag{3.24}
\end{equation}
Se siamo in un conduttore il lavoro del campo elettrico va in energia interna
(effetto Joule), altrimenti in energia cinetica delle cariche e, forse (ne
riparleremo), anche in qualche altra forma di energia (radiazione).\\
La (3.24) \`e la forma locale della (3.22); se siamo in un conduttore vale la
(3.17), e quindi
\begin{equation}
  w_j=\varrho_c\vec J^{\,2}=\sigma_c\vec E^{\,2}
  \tag{3.25}
\end{equation}
che costituisce la forma locale della (3.22').

\subsection{Un modello classico per la conduzione}

Per capire sia l'origine della legge di Ohm che i suoi limiti di validit\`a,
occorre capire in termini un po' pi\`u quantitativi il meccanismo della
conduzione. Consideriamo un conduttore omogeneo e di sezione costante; dalla
(3.17) e dalla (3.3) si ha $\vec J=\rho\vec u=\sigma_c\vec E$, da cui segue che
in presenza di un campo elettrico i portatori di carica si muovono con
velocit\`a (media) costante in direzione del campo, proporzionale al campo
stesso. Anche se il comportamento dei portatori di carica in un conduttore (che
pu\`o essere un metallo, ma anche un liquido od un gas ionizzato) \`e molto
complicato e soprattutto richiederebbe una trattazione Quantistica, possiamo
capirne gli elementi essenziali in un semplice modello classico dovuto a Drude e
Lorentz. Per comodit\`a ci riferiremo al caso in cui i portatori di carica sono
gli elettroni. L'idea fondamentale di questo modello \`e questa: tra un urto ed
il successivo con gli ioni del reticolo gli elettroni si muovono di moto
uniformemente accelerato sotto l'azione del campo elettrico $\vec E$; l'urto con
gli ioni altera in modo casuale la direzione della velocit\`a degli elettroni,
in modo tale che dopo un certo numero di urti, cio\`e \underline{in media} dopo
un certo tempo $\tau$, la direzione del moto dell'elettrone \`e completamente
scorrelata, cio\`e indipendente, dalla direzione iniziale. Il significato
preciso di questa ipotesi, di natura statistica, \`e questo: se immaginiamo un
gran numero di elettroni che inizialmente si muovono tutti in una certa
direzione, dopo un tempo $\tau$, in seguito agli urti, le loro direzioni di moto
sono uniformemente distribuite in tutte le direzioni. Per semplicit\`a
assumeremo che sia sufficiente un solo urto perch\'e si realizzi questa
``perdita della memoria'' delle condizioni iniziali, cio\`e assumiamo che $\tau$
sia uguale al tempo medio fra due urti. Sia $\vec v_i(t)$ la velocit\`a
dell'i--esimo elettrone all'istante $t$ e $\vec v_i^{\,c}$ la sua velocit\`a
subito dopo l'ultima collisione, avvenuta al tempo $t_i<t$; si ha
\begin{equation}
  \vec v_i(t)=\vec v_i^{\,c}-\frac em\vec E\cdot(t-t_i).
  \tag{3.26}
\end{equation}
Indichiamo con $\vec u(t)$ la velocit\`a media degli elettroni all'istante $t$
cio\`e, se $N$ \`e il numero di elettroni,
\begin{equation}
  \vec u(t)\defeq\frac1N\sum_{i=1}^N\vec v_i(t)
  \tag{3.27}
\end{equation}
e prendiamo la media di ambo i membri della (3.26):
\begin{equation}
  \vec u(t)=\frac1N\sum_{i=1}^N\vec v_i^{\,c}
  -\frac em\vec E\,\frac1N\sum_{i=1}^N(t-t_i).
  \tag{3.28}
\end{equation}
Il primo termine a secondo membro della (3.28) \`e la velocit\`a (vettoriale)
media degli elettroni subito dopo un urto, che \`e nulla per la nostra ipotesi
statistica di distribuzione uniforme in tutte le direzioni; il secondo termine
contiene il tempo medio trascorso dall'ultimo urto; esso \`e uguale al tempo
medio che deve trascorrere fino al prossimo urto, ed entrambi sono uguali a
$\tau$, il tempo medio fra due urti (e non $\tau/2\,$!): per convincersene si
pensi al gioco del Lotto, in cui un dato numero su una certa ruota esce in media
ogni 18 settimane, e la probabilit\`a che esca \`e sempre 1/18, anche se \`e
uscito la settimana precedente o se \`e in ritardo di 100 settimane (inutile
giocare i numeri in ritardo!), quindi il tempo medio di attesa \`e sempre di 18
settimane, a partire da un momento qualsiasi. Quindi la (3.28) diventa
\begin{equation}
  \vec u=-\frac{e\tau}{m}\vec E
  \tag{3.29}
\end{equation}
(il rapporto costante fra la velocit\`a ed il campo, cio\`e $e\tau/m$, \`e
quello che quando abbiamo introdotto i conduttori abbiamo chiamato la
\textbf{mobilit\`a} degli elettroni).\\
Il risultato espresso dalla (3.29) ci ricorda il moto di una particella in un
mezzo viscoso soggetta ad una forza costante, in cui la velocit\`a limite \`e
proporzionale alla forza: il modello che abbiamo appena discusso ci fornisce un
possibile meccanismo microscopico di ci\`o che macroscopicamente appare come
attrito viscoso.\\
Sia ora $n_c$ il numero di elettroni di conduzione per unit\`a di volume, allora
$-en_c$ \`e la densit\`a di carica e $-en_c\vec u$ \ la densit\`a di corrente;
quindi dalla (3.29) e dalla (3.17) si ha
\begin{equation}
  \sigma_c=\frac{e^2n_c\tau}{m}\,\cdot
  \tag{3.30}
\end{equation}
La relazione lineare fra campo elettrico e corrente, contenuta nella (3.29), \`e
la legge di Ohm, i cui limiti di validit\`a sono quindi quelli del modello. Le
ipotesi fatte per arrivare alla (3.29) sono molto generali, e difficilmente
confutabili (avere supposto che $\tau$ coincida con il tempo medio fra due urti
\`e irrilevante), se non per il fatto di avere utilizzato i concetti della
fisica classica, cio\`e avere trattato gli elettroni come delle palline in un
flipper; quindi dobbiamo chiederci in quali condizioni la conducibilit\`a data
dalla (3.30) pu\`o risultare dipendente dal campo elettrico applicato. Nella
(3.30) solo $\tau$ ed $n_c$ possono eventualmente dipendere da $E$: in effetti
se il campo elettrico \`e troppo intenso il tempo medio fra due urti pu\`o
venire modificato sensibilmente; ci\`o accade quando l'incremento di velocit\`a
degli elettroni fra un urto e l'altro \`e paragonabile alla velocit\`a termica
degli elettroni cio\`e, in termini di energia, quando
\begin{equation}
  eEl\approx kT
  \tag{3.31}
\end{equation}
dove $l$ \`e il \textbf{cammino libero medio} degli elettroni, cio\`e la
distanza media percorsa fra due urti, e $(3/2)kT$ l'energia cinetica media degli
elettroni: a temperatura ambiente $kT=0.025$\,eV, \ $l\approx10^{-8}$\,cm in un
metallo o in un liquido, mentre $l\approx10^{-6}$\,cm in un gas; in quest'ultimo
caso (che \`e il caso pi\`u sfavorevole) la validit\`a della legge di Ohm
richiede
\begin{equation}
  E\ll\frac{kT}{el}\approx0.025\cdot10^6\,\mathrm{V/cm}\simeq25\,\mathrm{kV/cm}.
  \tag{3.32}
\end{equation}
Un campo molto intenso pu\`o produrre (specialmente in un gas in cui il cammino
libero medio \`e abbastanza grande) un altro effetto: se l'energia cinetica che
gli elettroni (o gli ioni portatori di carica) acquistano fra un urto e l'altro
\`e abbastanza grande, nell'urto contro altri ioni o atomi possono ionizzarli,
quindi aumenta il numero di portatori di carica, aumenta di conseguenza la
probabilit\`a di altri urti con conseguente ionizzazione, e cos\`i via: si
instaura un processo catastrofico risultante in una scarica all'interno del
conduttore (il gas). \`E quello che succede in una lampada al neon, e in questi
casi la legge di Ohm non \`e pi\`u applicabile ($n_c$ dipende da $E$).

Un'altra situazione in cui il modello di Drude -- Lorentz cessa di essere
applicabile si presenta quando il campo elettrico \`e rapidamente variabile: per
scrivere la (3.26) abbiamo supposto il campo elettrico (praticamente)
indipendente dal tempo fra un urto e l'altro; \`e chiaro che, per esempio nel
caso di campi elettrici variabili nel tempo sinusoidalmente, se il periodo \`e
confrontabile con $\tau$ (o minore), il meccanismo della conduzione sar\`a
radicalmente modificato. Inoltre, come vedremo pi\`u avanti, in questi casi si
ha anche un altro effetto: la presenza di un campo elettrico autoindotto che
altera la distribuzione della corrente all'interno del conduttore (effetto
pelle), e quindi la resistenza del conduttore stesso, che risulta quindi
dipendere dalla frequenza $\omega$ del campo applicato.

Qualitativamente il modello di Lorentz -- Drude per la conducibilit\`a ci
permette di capire l'aumento della resistenza dei conduttori metallici con la
temperatura, in quanto un aumento della temperatura comporta una diminuzione del
tempo medio fra due urti e quindi, per la (3.30), una diminuzione della
conducibilit\`a. Non possiamo per\`o andare molto lontano con questo argomento,
che certamente non \`e in grado di spiegarci il fenomeno della
\textbf{superconduttivit\`a} a basse temperature: la conducibilit\`a di molti
metalli subisce un improvviso ed enorme aumento ad una data temperatura (detta
``temperatura critica $T_c$''), dipendente dal tipo di metallo, dell'ordine di
qualche Kelvin ($T_c\simeq20$\,K per il rame): al di sotto di questa temperatura
critica la resistenza del conduttore diventa praticamente nulla. Si tratta di
una vera e propria transizione di fase, come la solidificazione di un liquido, e
la fisica classica non \`e in grado di spiegare le transizioni di fase.

Un altro caso in cui si manifesta l'inadeguatezza del modello classico \`e
quello del Na (sodio): la conducibilit\`a del sodio a temperatura ambiente
risulta essere $\sigma_c=1.9\times10^{17}$\,ues, e la densit\`a di portatori di
carica (uno per atomo in questo caso) \`e
$n_c=2.5\times10^{22}\,\mathrm{cm^{-3}}$ per cui dalla (3.30) risulta
\ $\tau\simeq3\times10^{-14}$\,s \ che corrisponde, alla velocit\`a di
agitazione termica di circa $10^7$\,cm/s, ad un cammino libero medio di una
trentina di passi reticolari! Come fa una pallina in un flipper denso di
ostacoli a superarne impunemente e regolarmente una trentina, senza incontrarne
nessuno? Sar\`a la Meccanica Quantistica a fornire la spiegazione, sostituendo
alla descrizione degli elettroni come palline, una descrizione degli elettroni
come onde, e le onde, grazie ai fenomeni di interferenza a cui danno luogo, in
una struttura ordinata come quella di un cristallo si comportano in maniera
ordinata.

\subsection{La forza elettromotrice}

Se consideriamo una parte di un circuito percorso da corrente, come abbiamo
fatto quando abbiamo introdotto la legge di Ohm, o quando abbiamo discusso
l'effetto Joule, la corrente che vi circola \`e dovuta al campo elettrico nel
conduttore o, equivalentemente, alla differenza di potenziale agli estremi della
parte di circuito considerata. Ma se consideriamo l'intero circuito, percorso da
una corrente stazionaria, non basta il campo elettrico per spiegare la corrente
che vi circola: infatti un circuito percorso da una corrente stazionaria, o
anche da una corrente lentamente variabile ma a divergenza nulla, deve
necessariamente essere un circuito chiuso: se vale la (3.11) i tubi di flusso di
$\vec J$ non possono interrompersi, quindi le cariche che si muovono in un
circuito elettrico compiono traiettorie chiuse, ed inoltre, a meno che il
circuito non sia superconduttore, dissipano energia per effetto Joule. Ma il
campo elettrico (statico) \`e conservativo, e quindi il lavoro che fa sulle
cariche lungo un percorso chiuso \`e nullo, pertanto una corrente in un circuito
richiede la presenza di forze non conservative \underline{motrici} agenti sulle
cariche, cio\`e che fanno lavoro positivo sulle cariche quando queste percorrono
l'intero circuito. Un dispositivo in grado di mettere in gioco queste forze \`e
detto un \textbf{generatore di forza elettromotrice}: le normali pile o la
batteria di un'automobile sono generatori di forza elettromotrice.\\
Si definisce \textbf{forza elettromotrice in un circuito} (f.e.m.),
\underline{non} una forza, \underline{ma} il lavoro fatto dal generatore
sull'unit\`a di carica quando questa percorre l'intero circuito; la f.e.m.,
normalmente indicata con $\mathcal E$, \`e quindi la circuitazione della forza
sull'unit\`a di carica, e il suo segno dipende dalla scelta del verso di
percorrenza del circuito:
\begin{equation}
  \mathcal E=\frac1q\oint\vec F_q\cdot d\vec s
  \tag{3.33}
\end{equation}
dove $\vec F_q$ \`e la forza totale agente sulla carica $q$ (la forza dovuta al
campo elettrico possiamo a piacere includerla o no in $F_q$, in quanto la
circuitazione di $\vec E$ \`e nulla). Il circuito sul quale \`e definita la
f.e.m.\ in generale \`e una \underline{linea chiusa}, non necessariamente un
filo conduttore. La f.e.m.\ ha le dimensioni di un potenziale, e si misura nelle
stesse unit\`a: in volt nel SI, in statvolt nel sistema di Gauss. Una pila da
9\,V \`e quindi un generatore in grado di fare un lavoro di 9\,J se la carica di
1\,C percorre un circuito come quello di figura 3.3, nel quale abbiamo riportato
il simbolo usato per la pila.
\fig{fig3_3}{3.3}
La definizione di f.e.m.\ data dalla (3.33) \`e generale e sar\`a
particolarmente significativa nel caso dei campi variabili (legge di
Faraday--Neumann), in cui le forze non conservative agiscono lungo tutto il
circuito; nel caso di una pila o di una batteria le forze non conservative sono
localizzate nel generatore e vanno da quello che abbiamo indicato come polo
negativo al polo positivo, quindi la f.e.m.\ $\mathcal E$ \`e l'integrale (3.33)
solo all'interno del generatore, limitato alle sole forze non conservative
$\vec F_q^{\,\mathrm{n.c.}}$; se nel circuito di figura 3.3 il verso di
percorrenza \`e quello antiorario,
\begin{equation}
  \mathcal E=\frac1q\int_-^+\vec F_q^{\,\mathrm{n.c.}}\cdot d\vec s.
  \tag{3.34}
\end{equation}
Se il circuito \`e aperto, cio\`e il generatore non \`e collegato al resto del
circuito, siccome all'interno del generatore le cariche sono in equilibrio, \`e
presente un campo elettrico che esercita sulle cariche una forza uguale e
opposta a $\vec F_q^{\,\mathrm{n.c.}}$, quindi $\mathcal E$ coincide con la
differenza di potenziale (d.d.p.) fra il polo positivo e quello negativo, a
circuito aperto. Ma \`e chiaro che la f.e.m.\ non va confusa con la differenza
di potenziale: sono due concetti diversi, il primo ha significato per le forze
non conservative, il secondo \`e definito solo per i campi di forze
conservativi.
\fig{fig3_4}{3.4}
Quando il circuito \`e chiuso, all'esterno del generatore le cariche circolano
dal polo positivo a quello negativo sotto l'azione del campo elettrico, mentre
all'interno del generatore le cariche (di solito non le stesse che circolano
all'esterno) si muovono dal polo negativo a quello positivo, cio\`e, grazie alla
presenza delle forze non conservative, contro il campo elettrico.\\
All'interno del generatore il moto delle cariche comporta una pi\`u o meno
grande dissipazione di energia per effetto Joule: siccome $I$ \`e la quantit\`a
di carica che circola nell'unit\`a di tempo, dal bilancio energetico riferito al
moto delle cariche dentro il generatore e all'unit\`a di tempo
\begin{equation}
  I\mathcal E+I\int_-^+\vec E\cdot d\vec s\equiv I(\mathcal E-V)=W_J>0
  \tag{3.35}
\end{equation}
(dove $V$ \`e la d.d.p.\ fra il polo positivo e quello negativo), si vede che
quando il circuito \`e chiuso \ $V<\mathcal E$: la differenza $\mathcal E-V$ \`e
detta ``caduta di potenziale'' all'interno del generatore, e dipende dalla
corrente $I$, difatti \`e nulla per $I=0$. Se assumiamo che il comportamento del
mezzo in cui si muovono le cariche all'interno del generatore sia ohmico, cio\`e
schematizzabile con una resistenza $r$, detta \textbf{resistenza interna} del
generatore, dalla (3.35) e (3.22') si ha
\begin{equation}
  I(\mathcal E-V)=I^2r\quad\ \Rightarrow
  \quad\ \mathcal E-V=Ir\quad\ \Leftrightarrow\quad\ V=\mathcal E-rI
  \tag{3.36}
\end{equation}
e quindi in questa ipotesi la caduta di potenziale \`e proporzionale alla
corrente erogata. In una pila nuova la resistenza interna \`e di qualche decimo
di ohm, e cresce con l'uso fino a raggiungere il valore di qualche ohm (valori
solo indicativi, in quanto dipendono molto dal tipo di pila).\\
Nella figura 3.5 abbiamo riportato il \underline{circuito equivalente} ad un
generatore di f.e.m.: un generatore ``ideale'' (cio\`e con resistenza interna
nulla) con in serie una resistenza uguale alla resistenza interna; quando il
generatore eroga una corrente $I$, la d.d.p.\ fra i poli $+$ e $-$ del
generatore \`e data dalla (3.36).
\fig{fig3_5}{3.5}
L'equazione di un circuito come quello della figura 3.3, cio\`e l'equazione che
lega la f.e.m.\ $\mathcal E$ del generatore (di resistenza interna $r$) e la
corrente nel circuito (di resistenza $R$), per la legge di Ohm e la (3.36) \`e
\begin{equation}
  V\equiv\mathcal E-rI=RI\quad\ \Rightarrow\quad\ \mathcal E=(R+r)I.
  \tag{3.37}
\end{equation}
Oppure: la potenza erogata dal generatore, che per definizione di f.e.m.\ \`e
$W=\mathcal EI$, deve eguagliare la potenza dissipata sia all'interno del
generatore che nella resistenza
esterna, e quindi per la (3.22') e le (3.35), (3.36)
\begin{equation}
  \mathcal EI=I^2r+I^2R
  \tag{3.38}
\end{equation}
da cui dividendo per $I$ si riottiene la (3.37).

Ma qual \`e l'origine delle forze non conservative all'interno del generatore o,
equivalentemente, chi fornisce l'energia? La risposta a questa domanda varia a
seconda del generatore: nel generatore di Van de Graaff (1931) il lavoro \`e
fatto da chi (o cosa) fa girare la cinghia che trasporta le cariche che vengono
poi depositate su una grossa sfera cava, contro il campo elettrico prodotto
dalla sfera carica; in una pila chimica, che \`e un sistema termodinamico in
ambiente a temperatura e pressione costanti, il lavoro eguaglia (se le
trasformazioni sono reversibili) la variazione dell'energia libera di Gibbs
cambiata di segno: $\mathcal L=-\Delta\mathcal G$, quindi c'\`e scambio di
calore e di lavoro con l'ambiente esterno, ma quasi tutto il lavoro
``elettrico'' proviene dalla variazione dell'energia interna (energia chimica)
del sistema.\\
Un esempio di generatore di f.e.m.\ molto istruttivo, anche se certamente non di
interesse commerciale come le pile, \`e costituito dalla \textbf{pila a
gravit\`a}: un tubo di vetro verticale contiene una soluzione di Ag$^+$NO$_3^-$
e porta agli estremi due elettrodi di Ag; la forza di gravit\`a provoca un
aumento della concentrazione degli ioni Ag$^+$ lungo il tubo dall'alto verso il
basso; di conseguenza i due elettrodi non possono essere entrambi in equilibrio
con la soluzione: l'elettrodo superiore manda in soluzione ioni Ag$^+$, restando
carico negativamente, mentre nella parte inferiore del tubo, dove la
concentrazione \`e maggiore, ioni Ag$^+$ si depositano sull'elettrodo,
caricandolo positivamente. Quindi si instaura un campo elettrico dal basso verso
l'alto, necessario per l'equilibrio fra elettrodi e soluzione.\\
Quando il circuito viene chiuso, per effetto del campo elettrico gli elettroni
passano lungo il circuito esterno dall'elettrodo superiore, dove sono in
eccesso, all'elettrodo inferiore: per ogni elettrone che lascia l'elettrodo
superiore, per ripristinare l'equilibrio fra elettrodo e soluzione, uno ione
Ag$^+$ passa in soluzione, e per ogni elettrone che arriva all'elettrodo
inferiore uno ione viene neutralizzato e quindi un altro ione passa dalla
soluzione all'elettrodo; in definitiva, per ogni elettrone che percorre il
circuito esterno, un atomo di Ag ``cade'' dall'elettrodo superiore a quello
inferiore ed il lavoro fatto su di esso \`e $Mgh$, dove $M$ \`e la massa
dell'atomo di Ag. In questo caso la forza non conservativa motrice \`e la forza
di gravit\`a! Questo \`e possibile perch\'e l'intero circuito chiuso non \`e
percorso da una stessa carica: fuori gli elettroni (leggeri), dentro gli ioni
argento (pesanti). \`E proprio questo fatto che permette alla gravit\`a, un
campo di forze conservativo, di giocare il ruolo di forza motrice: siccome non
\`e la stessa massa che compie l'intero circuito chiuso, il lavoro fatto dalla
gravit\`a non \`e nullo. In effetti il sistema non torna allo stato di partenza:
poco a poco l'elettrodo superiore si consuma e aumenta la massa di quello
inferiore.
\fig{fig3_6}{3.6}
Quindi la f.e.m.\ della pila a gravit\`a \`e
\begin{equation}
  \mathcal E=\frac{Mgh}{e}\,\cdot
  \tag{3.39}
\end{equation}
Moltiplichiamo numeratore e denominatore a secondo membro della (3.39) per il
numero di Avogadro $N_\mathrm A$: siccome il peso atomico dell'Ag \`e 108, e
$N_\mathrm Ae=96500$\,C (costante di Faraday), per $h=1$\,m si trova
$\mathcal E=11\times10^{-6}$\,V.

\subsection{Carica e scarica di un condensatore}

Il circuito di figura 3.7 (circuito $RC$) \`e costituito da una pila di
f.e.m.\ $V_0$, una resistenza $R$ (che include l'eventuale resistenza interna
della pila), un condensatore di capacit\`a $C$, ed un interruttore I a due vie:
nella posizione A la pila \`e inserita nel circuito, che quindi \`e il circuito
di carica del condensatore; nella posizione B la pila \`e esclusa ed il
condensatore si scarica attraverso la resistenza $R$.\\
Portiamo l'interruttore nella posizione A; fissiamo arbitrariamente come verso
di percorrenza del circuito che definisce il segno della corrente quello
(orario) che va esternamente dal polo positivo della pila a quello negativo;
indichiamo con $V(t)$ la d.d.p.\ all'istante $t$ fra le armature del
condensatore secondo il verso stabilito ($V(t)$ \`e la d.d.p.\ fra l'armatura
superiore e quella inferiore), e scriviamo l'equazione del circuito eguagliando
la d.d.p.\ (o f.e.m.) $V_0$ con la somma delle d.d.p.\ ai capi dei vari elementi
del circuito (i conduttori che collegano i vari elementi del circuito di solito
hanno resistenza trascurabile, in caso contrario la loro resistenza viene
inglobata in $R$):
\fig{fig3_7}{3.7}
\begin{equation}
  V_0=RI(t)+V(t)\,;\qquad\quad V(t)\equiv\frac{Q(t)}{C}\,\cdot
  \tag{3.40}
\end{equation}
Poich\'e $Q(t)$ \`e la carica sull'armatura superiore del condensatore, si ha
$I(t)=\dot Q(t)=C\dot V(t)$, quindi possiamo prendere come incognita
indifferentemente $I(t)$, $Q(t)$ o $V(t)$: se prendiamo $V(t)$, l'equazione del
circuito \`e
\begin{equation}
  \dot V(t)=-\frac{1}{RC}V(t)+\frac{V_0}{RC}
  \tag{3.41}
\end{equation}
che \`e identica all'equazione per la velocit\`a di un corpo in un mezzo
viscoso, soggetto ad una forza costante: il prodotto $RC$ ha le dimensioni di un
tempo (indipendentemente dal sistema di unit\`a), ed \`e l'analogo del tempo di
smorzamento nel moto in un mezzo viscoso. Esso \`e detto \textbf{tempo di
scarica} (o di carica) del condensatore, e viene indicato con $\tau$: se
prendiamo, per esempio, $R=10^3\,\Omega$\,; $C=10$\,nF, si ottiene
$\tau=10^{-5}$\,s.\\
La soluzione generale della (3.41), che \`e un'equazione lineare non omogenea,
\`e
\begin{equation}
  V(t)=V_0+Ae^{-t/\tau}.
  \tag{3.42}
\end{equation}
Con la condizione iniziale $V(0)=0$ (il condensatore inizialmente \`e scarico),
si ha
\begin{equation}
  V(t)=V_0\big(1-e^{-t/\tau}\big)\,;
  \qquad\ I(t)=C\dot V(t)=\frac{V_0}{R}e^{-t/\tau}.
  \tag{3.43}
\end{equation}
Nella figura 3.8 abbiamo riportato sullo stesso grafico sia l'andamento di
$V(t)$ che quello di $RI(t)$.
\fig{fig3_8}{3.8}
Ci sono un paio di cose da discutere a proposito del modo in cui abbiamo
ottenuto la (3.43): in primo luogo abbiamo applicato la legge di Ohm in un caso
in cui la corrente non \`e stazionaria; questo \`e lecito purch\'e la corrente
sia una corrente lentamente variabile, che in questo caso significa che $\tau$
deve essere abbastanza grande. Grande rispetto a cosa? Anticipando cose che
giustificheremo in seguito, $\tau$ deve essere grande rispetto al tempo che i
segnali elettromagnetici (che viaggiano alla velocit\`a della luce) impiegano ad
attraversare il circuito: se il circuito ha le dimensioni di 30\,cm questo tempo
\`e di $10^{-9}$\,s mentre, come abbiamo visto, con $R=10^3\,\Omega$\,,
\ $C=10$\,nF, \ \ $\tau=10^{-5}$\,s.\\
Si noti che anche se $I(t)$ \`e lentamente variabile, $\vec J$ non \`e a
divergenza nulla perch\'e il circuito (cio\`e il tubo di flusso) si interrompe
sulle armature del condensatore: e infatti sull'armatura superiore si ha un
accumulo di carica, e su quella inferiore una diminuzione, in accordo con la
(3.10).\\
La seconda cosa che vogliamo discutere \`e la discontinuit\`a in $I(t)$: la
corrente \`e nulla finch\'e non si chiude l'interruttore, poi improvvisamente
passa al valore $V_0/R$, e questo \`e un po' sospetto in quanto,
schematizzazioni a parte, di solito le grandezze fisiche variano con
continuit\`a. Anche qui interviene il fatto che la corrente non \`e stazionaria;
quando la corrente in un circuito dipende dal tempo, diventa importante un'altra
grandezza caratteristica del circuito, oltre alla resistenza e la capacit\`a:
l'induttanza, il cui effetto, pi\`u o meno grande a seconda di quanto
rapidamente varia la corrente, \`e quello di introdurre una specie di inerzia
alle variazioni della corrente con il risultato, mostrato nella figura 3.9, di
rendere continuo il passaggio della corrente dal valore 0 al valore (quasi)
$V_0/R$.
\fig{fig3_9}{3.9}
Dopo un tempo $t_0\gg\tau$ il condensatore \`e (praticamente) carico al
potenziale $V_0$; a questo punto portiamo l'interruttore nella posizione B. La
nuova equazione del circuito \`e
\begin{equation}
  \dot V(t)=-\frac{1}{RC}V(t)
  \tag{3.44}
\end{equation}
cio\`e la (3.41) con $V_0=0$. Presa $t_0$ come nuova origine dei tempi, la
soluzione della (3.44) con la condizione iniziale \ $V(0)=V_0$ \ \`e
\begin{equation}
  V(t)=V_0e^{-t/\tau}\,;\qquad\ I(t)=C\dot V(t)=-\frac{V_0}{R}e^{-t/\tau}
  \tag{3.45}
\end{equation}
cio\`e corrente e potenziale si portano esponenzialmente a zero con il solito
tempo caratteristico $\tau$. Il segno $-$ nell'espressione di $I(t)$ \`e dovuto
al fatto che ora la corrente circola in senso antiorario, e noi non abbiamo
cambiato la convenzione sul verso delle correnti positive (potevamo farlo).
Anche in questo caso la $I(t)$ presenta una discontinuit\`a, che in realt\`a non
esiste.

Calcoliamo l'energia dissipata per effetto Joule nella resistenza durante la
scarica del condensatore:
\begin{equation}
  \begin{split}
    E_\mathrm J&=\int_0^\infty W_\mathrm J(t)\,dt
    =\int_0^\infty\frac{V(t)^2}{R}\,dt
    =\frac{V_0^2}{R}\int_0^\infty e^{-2t/\tau}\,dt\\
    &=\frac{V_0^2\tau}{2R}=\frac12CV_0^2
  \end{split}
  \tag{3.46}
\end{equation}
che coincide, come dovevamo aspettarci, con l'energia immagazzinata nel
condensatore nella fase di carica, data dalla (2.32).\\
E quant'\`e l'energia dissipata nella resistenza nella fase di carica del
condensatore? Il conto si pu\`o fare nello stesso modo e non presenta
difficolt\`a, ma in effetti non c'\`e bisogno di fare nessun conto: nella fase
di carica la pila fa circolare la carica $Q=CV_0$, quindi (definizione di
f.e.m.) compie il lavoro $QV_0=CV_0^2$; met\`a di questo lavoro
($\frac12CV_0^2$) va in energia del condensatore, quindi l'altra met\`a viene
dissipata nella resistenza.

L'importanza dei condensatori nella tecnica elettronica risiede sostanzialmente
nel fatto che consente di introdurre nell'equazione di un circuito un
\underline{tempo} caratteristico del circuito, cio\`e la possibilit\`a di
correnti variabili grazie alle propriet\`a intrinseche del circuito:
l'alternativa sarebbe quella di fare intervenire il tempo dall'esterno, per
esempio tramite un generatore di f.e.m.\ variabile.
